

%Let \(\Delta_X = \Delta(x_1,...,x_n) := \max \|x_i-x_j\| / \min \|x_i-x_j\|\) denote the aspect ratio of a point set in Euclidean space. For \(X \in \mathbb{R}^{n\times d}\) a matrix of points, let \(D_X \in \mathbb{R}^{n\times n}\) denote its corresponding interpoint distance matrix. 

%The problem of producing a t-SNE plot proceeds as follows: 
%\paragraph{Basic Notation}

%\paragraph{Notation} Unless otherwise specified \(\|\cdot\|\) refers to the Euclidean norm; big-O notation is with respect to \(n\), the number of points; 

%\paragraph{t-SNE} 
Given an input dataset \(X = \{x_1,\ldots, x_n\}\subset \mathbb{R}^{D}\), the goal of t-SNE is to come up with an embedding \(Y = \{y_1,\ldots,y_n\}\subset \mathbb{R}^{d}\) (where \(d\ll D\), typically \(d=2\))\footnote{With a slight abuse of subset notation, duplicates are allowed in $X$ and $Y$.} that approximately maintains the neighborhood structure in $X$. t-SNE accomplishes this by assigning affinities to input data points which encode how likely an input point is to be a neighbor to a given point. The goal then is to find a configuration of the embedded points $Y$ that induces a similar neighborhood affinity. 
Specifically, let \(P = P(X) \in \mathbb{R}_{\geq 0}^{n\times n}\) and \(Q = Q(Y) \in \mathbb{R}_{\geq 0}^{n\times n}\) be the input and embedded
\textit{affinity matrices} describing the pairwise neighborhood similarities in the input and output, respectively. t-SNE constructs \(P\) 
by first computing neighborhood affinities for each point $i$ defined as (for any $i\neq j$)
%based on a symmetrized Gaussian kernel with varying neighborhood sizes \(\{\sigma_i: i\in [n]\}\) for the input affinities. For any input data point $i$, define the neighborhood affinities (with respect to $i$) as
%conditional Let us first define for distinct \(i,j \in [n]\) the \textit{conditional Gaussian kernel} 
\begin{equation}
    P_{j|i}(X; \sigma_i) := \frac{\exp(-\lVert x_j-x_i\rVert^2/(2\sigma_i^2))}{\sum_{k\neq i} \exp(-\lVert x_k -x_i\rVert^2/(2\sigma_i^2))} 
    \hspace{0.5in} P_{i|i}(X; \sigma_i):=0
    \label{eq:tsne_p}
\end{equation}
where $\sigma_i \geq 0$ encodes the (point-dependent) neighborhood scaling\footnote{If $\sigma_i=0$, define $P_{j|i}(X;0) := \lim_{\sigma_i \to 0} P_{j|i}(X;\sigma_i)$.%If $\sigma_i$ is chosen as $0$, then one takes the limit of Eq.\ (\ref{eq:tsne_p}).
}. When $X$ and $\sigma_i^*$ are clear from context, we will often drop it from the notation. It is worth noting that $P_{\cdot|i}$ is a valid probability distribution over $[n]$.
The matrix \(P\) is then constructed based on a crucial parameter called the perplexity (denoted as \(\rho\) and taking values in %\footnote{Notably, \citet{jeong2024convergence} established that the procedure is only well-defined for \(\rho\in [1,n-1]\)}
\([1,n-1]\)), as follows: 
\begin{itemize}
    \item[(1)] For each $i \in [n]$, select the (unique, see Lemma \ref{lem:unique_neighborhood}) neighborhood scale \(\sigma_i^*\) that minimizes the gap between the entropy of \(P_{\cdot | i}(X;\sigma^*_i)\) and \(\log_2 \rho\). %\todo{check}
    %For each $i$, select the neighborhood scale \(\sigma_i^*\) such that the entropy of the neighborhood distribution \(P_{\cdot | i}(X;\sigma^*_i)\) is \(\log \rho\). \textcolor{red}{    cc }%\todo{$P|_i$?} 
    \item[(2)] Define \(P =  [P_{ij}]_{i,j\in [n]}\) where \(P_{ij} := \frac{1}{2n}(P_{i|j}(\sigma^*_j) + P_{j|i}(\sigma^*_i))\).% if \(i\neq j\) and zero otherwise.
\end{itemize}
To avoid the so-called \emph{crowding problem} (see \cite{van2008visualizing} for details), the output affinity matrix \(Q\) is computed based on a t-distribution. Specifically, for \(i\neq j\)
\begin{equation}
    Q_{ij}(Y) := \frac{(1 + \|y_i -y_j\|^2)^{-1}}{\sum_{k,l ; k\neq l} (1 + \|y_k-y_l\|^2)^{-1}} \hspace{0.5in}
    Q_{ii}(Y) :=0.
\end{equation}
As indicated before, the objective then is to minimize the gap between the input and output affinities $P$ and $Q$. This is accomplished by minimizing relative entropy (KL-divergence) between the \(P\) and \(Q\) affinities (viewed as probability distributions).  \[\textup{minimize}_Y \; \mathcal{L}_X(Y):= \KL(P(X)\|Q(Y)) = \sum_{\substack{i,j\\ i\neq j}} P_{ij}(X) \log\Big( \frac{P_{ij}(X) }  {Q_{ij}(Y)}\Big).\]
%\todo[inline]{L(Y) needs to have a reference to $X$}
This highly non-convex objective is usually optimized by initializing at a good starting point via an \emph{early exaggeration phase}, followed by performing standard gradient descent methods and returning an embedding $Y$ that corresponds to a local minimum of the objective. Our central task is to study the nature of the these (local minimum) embeddings returned by t-SNE and their relation to the space of input datasets.
%We can now define the main object of interest in this paper: optimal t-SNE embeddings.
\begin{definition}
    For an $n$-point dataset\footnote{Without loss of generality, we shall assume that the input dimension $D=n-1$.} $X \subset \mathbb{R}^{n-1}$ and perplexity parameter $\rho \in [1,n-1]$, define%\footnote{\(\TSNE_{\rho}(X)\neq \varnothing\), as the trivial embedding where all points are mapped to the origin is stationary.}
    $${\TSNE}_{\rho}(X) := \{ Y \subset \mathbb{R}^d : \nabla_Y \mathcal{L}_X(Y) = 0 \} $$
    as the set of outputs $Y\subset \mathbb{R}^d$ that are stationary to the t-SNE objective on a given input $X$.
    Furthermore, for a set of $n$-point datasets $\mathcal{X}_n$, we define:
    $${\TSNE}_{\rho}(\mathcal{X}_n) := \bigcup_{X \in \mathcal{X}_n}{\TSNE}_{\rho}(X).$$
    If $\mathcal{X}_n$ is the set of \emph{all} $n$-point datasets, we denote ${\TSNE}_\rho(\mathcal{X}_n)$ as ${\IMTSNE}$ to indicate the entire image of the t-SNE map.
\end{definition}
All the supporting proofs for our formal statements can be found in the Appendix, and the code related to our experimental demonstrations is available on Github at \href{https://github.com/njbergam/tsne-exaggerates-clusters}{\texttt{https://github.com/njbergam/tsne-exaggerates-clusters}}. %\href{https://github.com/anon594/iclr26_submission8125}{\texttt{https://github.com/anon594/iclr26\_submission8125}}.

%\paragraph{Code} The code related to our experimental demonstrations is publicly available on \href{https://github.com/anon594/iclr26_submission8125/commits/main/}{GitHub}.

